\documentclass[11pt]{article}
\usepackage{amsmath,amssymb,amsthm,hyperref}
\newtheorem{theorem}{Theorem}
\newtheorem{corollary}{Corollary}
\newcommand{\E}{\mathbb E}
\title{A global lower bound for every endpoint odds coordinate}
\author{Asymptotic research note; independently reviewed}
\date{30 September 2026}
\begin{document}\maketitle
Let $p_1,\ldots,p_m:[K]\to[0,1]$ be nondecreasing, with positive row
weights $w_q$ summing to one, and suppose
\[
 \E_w\prod_{j\in S}p_j=\frac1{|S|+1}\quad(S\subseteq[m]).
\]
On rows where $A=\prod_jp_j>0$, write $r_j=(1-p_j)/p_j$ and
$\lambda=\sum_jr_j$.

\begin{theorem}
For $m\ge2$, every row $q$ with $A(q)>0$ and $L=\lambda(q)\ge32$
satisfies
\[
 r_i(q)\ge\frac{\min\{L,m-1\}}{32m}\qquad(1\le i\le m).
\tag{1}
\]
\end{theorem}
\begin{proof}
Fix $i$. If $r_i(q)>L/2$, the assertion is immediate. Otherwise put
$\ell=\lfloor\min\{L,m-1\}/16\rfloor$ and choose $J$ to consist of the
$\ell$ smallest odds among the $m-1$ indices other than $i$, evaluated
at row $q$. Then
\[
 \lambda_{-(J\cup\{i\})}(q)
 \ge\left(1-\frac\ell{m-1}\right)(L-r_i(q))
 \ge\frac{15L}{32}.
\tag{2}
\]
For a failure set $T$, define the probability measure
\[
 \nu_T=(m+1)\binom m{|T|}\,w
       \prod_{j\in T}(1-p_j)\prod_{j\notin T}p_j.
\]
The mixed moments imply this normalization. They also imply
\[
 \E_{\nu_T}\lambda_{-T}\le |T|+1.
\tag{3}
\]
Indeed for each $j\notin T$ the expectation of $r_j$ is at most
$(|T|+1)/(m-|T|)$, by the ratio of the probabilities of the specified
failure patterns $T\cup\{j\}$ and $T$. The inequality permits zero
$p_j$ rows, whose mass is counted only by the latter pattern.

Use (3) for $T=J\cup\{i\}$. By monotonicity and (2), every row strictly
before $q$ in the support of $\nu_T$ has remaining odds sum at least
$15L/32$. As $\ell+2\le L/16+2\le L/8$, Markov's inequality gives
\[
 \nu_T(\{s<q\})\le\frac{\ell+2}{15L/32}\le\frac4{15}.
\]
In particular the suffix $s\ge q$ has $\nu_T$ mass at least $1/2$.
All $p_j$ are positive on that suffix. Put $d=m-\ell$; there
\[
 \nu_T=\frac d{\ell+1}r_i\nu_J.
\]
Since $r_i(s)\le r_i(q)$ on the suffix, we obtain
\[
 \frac12\le\nu_T(\{s\ge q\})
 \le\frac d{\ell+1}r_i(q)\nu_J(\{s\ge q\})
 \le\frac d{\ell+1}r_i(q).
\]
Now $d\le m$ and $\ell+1>\min\{L,m-1\}/16$, proving (1).
\end{proof}

\begin{corollary}[Rows with zero success product]
For $m\ge2$, if $A(q)=0$, then every column at that row satisfies
\[
 p_i(q)\le32/33.
\]
\end{corollary}
\begin{proof}
The claim is immediate if $p_i(q)=0$. Otherwise order the other odds
using the value $+\infty$ for zero entries, put
$\ell=\lfloor(m-1)/16\rfloor$, and let $J$ be the $\ell$ smallest of
them. Since at least one of the other columns is zero and
$\ell<m-1$, at least one zero remains outside $T=J\cup\{i\}$.
Monotonicity implies that this column vanishes on every row up to $q$.
Consequently $\nu_T$ gives all its mass to the strict suffix $s>q$.
On that suffix $p_i>0$, so the identity
$\nu_T=(m-\ell)r_i\nu_J/(\ell+1)$ is valid even if other columns
vanish: both sides then vanish. Monotonicity of $r_i$ gives
\[
 1\le\frac{m-\ell}{\ell+1}r_i(q),\qquad
 r_i(q)\ge\frac{\ell+1}{m-\ell}
 >\frac{m-1}{16m}\ge\frac1{32}.
\]
This is equivalent to the assertion.
\end{proof}

The argument uses a number of conditioned failures that grows with the
row's odds sum, and chooses those failures among the smallest odds.
It does not require an upper balance theorem and remains valid for
$\lambda\gg m$. This is a structural comparison-tensor theorem, not
by itself a new lower bound on the number of rows.
\end{document}
